%\usepackage{xskak}
\usepackage{casiovn}
%\usepackage{etoolbox}%tính toán, công cụ
%
%\usepackage{graphicx}
%\usepackage{tabvar}
\renewcommand{\vec}[1]{\protect\overrightarrow{#1}}
\newenvironment{boxtb}{\begin{tomtat}}{\end{tomtat}}
\newenvironment{makerr}{\begin{noidung}{CHÚ Ý}}{\end{noidung}}
\def\viduminhhoa{\subsubsection{Ví dụ minh họa}}
\def\baitaptl{\subsubsection{Bài tập áp dụng}}
\def\baitapbs{\subsubsection{Bài tập bổ sung}}
\def\BTVD{\subsubsection{Bài tập vận dụng}}
%-----Vidu mới
%--------------------Ví dụ mới
\newenvironment{vidu}{\begin{boxvidu}\begin{myvidu}}{\end{myvidu}\end{boxvidu}}
%%%%%%%%%LVD
\def\baitoan{\faHandORight\ \textbf{Bài toán. }}
\def\phuongphap{\faHandODown\ \textbf{Phương pháp: }}

\newcommand{\boxde}[1]{%
	\setcounter{ex}{0} % đặt bộ đếm Câu về 1
	\setcounter{bt}{0}
	%	\addcontentsline{toc}{subsection}{#1~\thede} % đưa MT vào mục lục
	\begin{center}
		\begin{tikzpicture}[baseline=(A.base),scale=1]%
			\node[thick,scale=1,fill=teal!80,draw=teal,minimum width=2.5cm,minimum height=0.1cm,rounded corners=2mm] (A) {\bf\fontfamily{qag}\selectfont\color{white}#1};
		\end{tikzpicture}%
	\end{center}
}
%-----Đn lại lệnh lời giải
\AtBeginEnvironment{vidu}{
	\def\loigiaiEX{\makebox[\linewidth][c]{\rule[-0.2cm]{0pt}{0.6cm}\color{brandNavy}\damEX\faPencilSquareO\ Lời giải}}
}
\foreach \moitruong in {ex,vd,bt}{
	\AtBeginEnvironment{\moitruong}{
		\def\loigiaiEX{\makebox[\linewidth][c]{\rule[-0.2cm]{0pt}{0.6cm}\color{brandNavy}\damEX\faPencilSquareO\ Lời giải}}
	}
}
%%%%%%%%%%% cộng đa thức cột
\newcommand{\dropsign}[1]{\smash{\llap{\raisebox{-.5\normalbaselineskip}{$#1$\hspace{2\arraycolsep}}}}}

%-----------------
\newenvironment{luyentap}{\begin{boxvidu}\begin{myvidu}}{\end{myvidu}\end{boxvidu}}
\newenvironment{kttrongtam}{\begin{boxdn}}{\end{boxdn}}
\newenvironment{hd}{\begin{boxdn}}{\end{boxdn}}
\newenvironment{vandung}{\begin{bt}}{\end{bt}}
\newenvironment{trithuc}{\begin{boxdn}}{\end{boxdn}}
\renewenvironment{hd}{\begin{boxdl}}{\end{boxdl}}
\newenvironment{bttuongtu}{\begin{bt}}{\end{bt}}
\def\baitap{\subsection{BÀI TAPPPPPPPP}}

\def\BTTL{}